\section{Генераторы}
\label{sec:Generators}
\tikzsetfigurename{FLT_Generators_}

В предыдущем разделе мы рассмотрели распознаватели~--- формальные машины, принимающие или отвергающие слова и тем самым задающие язык <<извне>>, через ответ на вопрос о принадлежности.
Двойственным к такому подходу является подход генеративный: язык задаётся <<изнутри>>, через описание способа построения его слов.
Один из базовых способов задать генератор языка опирается на \emph{системы переписывания строк}, из которых, в дальнейшем, можно получить так называемые \emph{порождающие грамматики}.

\begin{definition}[Система переписывания]
    \label{def:RewritingSystem}
    \emph{Система переписывания (строк)}%
    \sidenote{
        Один из классических примеров систем переписывания строк~--- это машины Маркова, или алгорифмы Маркова~\cite{markov1954theory}.
        В отличие от нашего определения, в машинах Маркова правила упорядочены и применяются детерминированно: на каждом шаге выбирается первое применимое правило, и оно применяется к самому левому подходящему вхождению.
        Мы же рассматриваем недетерминированный вариант, в котором может быть применено любое правило к любому подходящему вхождению.
    }~---
    это пара $R = \langle \Sigma, P \rangle$, где
    \begin{itemize}
        \item $\Sigma$~--- алфавит,
        \item $P= \{p \mid p = w' \to w'', w' \in \Sigma^+, w'' \in \Sigma^*\}$~--- набор правил переписывания.
    \end{itemize}
\end{definition}

Один шаг работы системы состоит из замены любого вхождения любой из левых частей правил на соответствующую правую часть правила.
Иными словами, пусть есть слово $w = w_0 w_1 w_2$ и имеется правило $w_1 \to w_3 \in P$.
Тогда после переписывания по этому правилу будет получено слово $w' = w_0 w_3 w_2$.
Заметим, что выбор заменяемого вхождения, как и выбор применяемого правила, недетерминирован: может быть выбрано любое из вхождений и применено любое из доступных правил.
В качестве критерия остановки выберем невозможность применить ни одно из доступных правил.
Получившееся при этом слово будем считать результатом работы системы.
Зафиксировав стартовую строку и набор правил, можно достаточно естественным образом получить язык как множество всех слов, являющихся результатами работы системы при всевозможных последовательностях переписываний.

\begin{example}
    Рассмотрим систему переписывания $R = \langle \{a, b, c\}, P \rangle$ с правилами
    \[
        P = \{ab \to ba,\ b \to c\}
    \]
    и стартовой строкой $w_0 = ab$.
    К строке $ab$ применимы оба правила, причём правило $ab \to ba$~--- к единственному вхождению $ab$, а правило $b \to c$~--- к символу $b$ (единственному вхождению левой части $b$).
    Таким образом, возможны два варианта:
    \begin{itemize}
        \item применить $ab \to ba$: получим $ba$, после чего ни одно правило неприменимо~--- результат $ba$;
        \item применить $b \to c$: получим $ac$, после чего ни одно правило неприменимо~--- результат $ac$.
    \end{itemize}
    Других последовательностей переписываний нет, поэтому язык, порождаемый данной системой из стартовой строки $ab$, есть $L = \{ba, ac\}$.
\end{example}

Представленная система является, скорее, неформальной базой для создания более содержательных систем.
Далее, накладывая дополнительные ограничения на правила и алгоритм работы, мы будем получать различные содержательные классы языков.

\emph{Порождающая грамматика} (или просто \emph{грамматика})~--- это один из способов структурировать систему переписывания, выделив в алфавите две категории символов и зафиксировав стартовый символ.

\begin{definition}[Порождающая грамматика]
    \label{def:GenerativeGrammar}
    \emph{Порождающей грамматикой} называется четвёрка $G = \langle \Sigma, N, P, S \rangle$, где
    \begin{itemize}
        \item $\Sigma$~--- конечный \emph{терминальный алфавит} (символы, из которых строятся слова языка);
        \item $N$~--- конечный \emph{нетерминальный алфавит} (вспомогательные символы), причём $\Sigma \cap N = \varnothing$;
        \item $P$~--- конечное множество \emph{правил вывода} (продукций), каждое из которых имеет вид $\alpha \to \beta$, где $\alpha \in (\Sigma \cup N)^+$ и $\beta \in (\Sigma \cup N)^*$;
        \item $S \in N$~--- \emph{стартовый нетерминал}.
    \end{itemize}
    \emph{Язык, порождаемый грамматикой} $G$,~--- это множество всех слов над $\Sigma$, которые можно получить из стартового нетерминала $S$ последовательным применением правил из $P$:
    \[
        L(G) = \{w \in \Sigma^* \mid S \derives w\}.
    \]
\end{definition}

Нетрудно видеть, что грамматика является частным случаем системы переписывания: если отвлечься от деления символов на терминалы и нетерминалы, то $G$~--- это система переписывания над алфавитом $\Sigma \cup N$ с правилами $P$ и стартовой строкой $S$, в которой нас интересуют только те результаты, которые состоят исключительно из терминалов.
Разделение на терминалы и нетерминалы позволяет управлять <<незавершённостью>> промежуточных строк и гарантировать, что результирующие слова языка строятся из <<целевых>> символов.

При спецификации конкретных грамматик часто опускают явное перечисление терминалов и нетерминалов, указывая только множество правил.
При этом нетерминалы принято обозначать прописными латинскими буквами, терминалы~--- строчными, а стартовый нетерминал~--- буквой $S$. Последнее, хоть и часто используется, не обязательно. В некоторых случаях стартовым нетерминалом принято считать левую часть самого первого при записи правила.
Мы будем следовать этим соглашениям, если не указано иное.

Типичные примеры генераторов, встречающихся в повседневной жизни: генераторы паролей и почтовых адресов.
Эти инструменты по запросу генерируют строку, гарантированно%
\sidenote{Мы предполагаем, конечно, что генератор не содержит ошибок.}
обладающую необходимыми свойствами. Например, являющуюся корректным, с точки зрения заранее заданных правил, паролем.

Отметим, что подход к определению языков через системы переписывания не всегда удобен.
Например, не получается с их использованием естественным образом определить булевы языки~\sidecite{Okhotin:2003:BG:1758089.1758123}, поскольку булевы грамматики оперируют не только операцией вывода (переписывания), но и операциями конъюнкции и отрицания, не имеющими прямого аналога в терминах правил переписывания.
В следующих главах мы рассмотрим различные классы грамматик, получаемые наложением ограничений на вид правил, и соответствующие им классы языков.
